\section{引言}
\label{sec:introduction}
面向仓库的软件智能体正日益将代码检索、编辑、命令执行和反馈整合进同一工作流。SWE-bench 衡量完整系统解决真实仓库问题的能力，SWE-agent 和 OpenHands 则展示了语言模型周围的接口与执行环境的重要性 \cite{jimenez2024swebench,yang2024sweagent,wang2025openhands}。这些系统使另一个问题变得至关重要：当智能体依据先前观察到的结果采取行动时，什么将该结果与即将被修改的软件状态联系起来？

设想一个智能体正在修改会话刷新行为。它定位到正确的函数并取得源码，随后另一工作者更新了同一文件。此时，先前修订的测试结果仍可获取；第二个验证器针对更窄的目标报告通过；实现任务也被标记为完成。这些观察并不一定有误。错误发生在它们各自不同的含义被合并为一个信号，进而被理解为当前变更已经就绪之时。定位成功不能证明编辑所依据的源码仍是当前版本；执行过检查不能证明所有需求都已被覆盖；任务完成也不能证明验收通过。

本文从运行时边界研究这一协调问题。我们将\emph{证据}定义为针对工件或操作的带类型观察，其中记录了来源和适用条件。证据不再足以证明某项当前操作合理之后，仍可能作为历史记录发挥作用。因此，关键区别在于保留一项观察与将其纳入当前决策之间。只要智能体、人工编辑者和后台检查共同使用持续演化的软件状态，这一区别就至关重要。

本文提出 \wcode{}，这是一个已实现的 Rust 运行时，通过模型上下文协议（Model Context Protocol）提供接口。它维护软件意图声明、代码观察、验证记录和工作项状态。其组织原则是：操作所使用的证据必须具有明确的适用条件。读取操作返回原始源码与文件身份；修改操作将当前文件与已观察的身份进行核对；验证操作记录检查执行时绑定的修订与权威性；协调核对则保留生产者之间与当前操作相关的分歧。同一原则也决定接口应将哪些内容报告为缺失或不确定。

本文的研究问题是：\emph{共享的软件智能体运行时，如何在有界观察、并发修改与验证之间保持证据的适用性？} 本文通过以下三项贡献考察这一问题：
\begin{enumerate}
\item 给出明确的准入契约，将证据与代码和设计修订、目标范围、权威性及策略关联起来。该契约区分支持当前决策的证据、历史观察以及尚未解决的生产者冲突。
\item 实现一个运行时，将这些区别贯穿于上下文打包、受保护的单文件编辑、验证执行与复用以及协调核对。贯穿示例说明，为何必须在每个边界保留这些字段。
\item 提供可审计的诊断评估，分别测量源码交付、编辑输入就绪状态、运行控制和描述性开销。保留的可执行文件与源码归档标识主试验运行；通过移除谓词条件的分析，揭示被饱和检索分数掩盖的就绪判定错误。
\end{enumerate}

本文的贡献是围绕共享操作契约对已有机制进行组合。哈希、来源记录、架构一致性和访问控制均已有大量研究；第~\ref{sec:related} 节将本运行时置于这些研究脉络中。相应地，本文评估的范围也是明确限定的：评估不调用语言模型，只确认开发者编写的场景中的行为，不证明补丁成功率有所提高。这一范围使我们能够直接考察工具边界上的失败，而不将其归因于模型推理；对下游有效性的评估则留待独立实验。
